\section{はじめに}
\label{sec_intro}
大喜利は，出題された題目に対して機知に富んだ回答を返す言葉遊びであり，ユーモアの生成と理解を扱う研究の題材として用いられている\cite{ref_oogiri_go, ref_newyorker}．
大喜利投稿サイト「ボケて」\cite{ref_bokete}では，投稿者がお題となる画像を掲示し，他の利用者がその画像に対する面白い一言を投稿する．
以下では，お題となる画像をお題画像，投稿された一言をボケと呼ぶ．
ボケてから収集したデータベースには，約682万件のお題画像と約9,250万件のボケが含まれる．
筆者らはこれまでに，全てのお題画像に大喜利の形式や画像の媒体を表すラベルを付与した．
本稿では，その次の段階として，ボケそのものにラベルを付与する方法を扱う．

ボケの生成や面白さの判定を学習・評価する際には，ボケがどのような作りの文であるかを知る必要がある．
例えば，画像の人物のセリフとして書かれたボケと，画像に題名を付けたボケでは，生成の際に与えるべき条件が異なる．
また，テレビ番組や歌の知識を前提とするボケは，その知識を持たないモデルには面白さを判断できない．
このようなボケを区別せずに評価すると，知識の不足による誤りとユーモアの理解の誤りを分けられない．
さらに，カルタのお題で札のひらがなから始まらないボケや，お題と無関係な書き込みは，生成の学習データから除くべきである．
したがって，ボケの作りを表すラベルを全てのボケに付与できれば，学習データの絞り込み，生成の条件付け，判定の誤りの分析に利用できる．

しかし，ボケのラベルには，お題画像のラベルにない2つの難しさがある．
第一に，何をラベルとすべきかが自明ではない．
ボケの区分には，笑いの手法や面白さのように主観的なものも多く，人でも判断が揃いにくい．
第二に，対象が約9,250万件と多い．
視覚言語モデル (Vision-Language Model，VLM) は画像と指示文から分類を行えるが，大規模なモデルを全件に適用すると，単一の GPU で数百日を要する．
本研究では，有料の API を用いず，単一の GPU で全件を現実的な時間で処理できる方法を目標とした．

本稿は，ラベルの設計から自動判定の方法の確立までの試行錯誤の経緯を，実験の結果とともに報告する．
主な内容は次のとおりである．
\begin{itemize}
    \item 実例の読解から，お題の形式が課す制約の順守，発話の型，画像外の知識の3つの観点を定め，定義書を3回改訂した経緯と，その判断の根拠を示す．
    \item 検証データ620件を作成し，同一の判断者による付け直しの一致度を測ることで，自動判定の精度の実質的な上限を示す．
    \item 制約の順守を，規則と同じお題の他のボケとの類似度で判定できることを示す．
    \item VLM の教師への指示，教師ラベルによる文の埋め込みの分類器，小さな VLM の微調整を比較し，教師ラベルと少量の目視ラベルによる二段階の学習が最も有効であることを示す．
    \item 画像外の知識の判定を改善するための試みと，その限界を示す．
    \item 同じお題のボケで計算を共有する推論の高速化を示す．
\end{itemize}

本稿の構成は次のとおりである．
\ref{sec_data}章で対象のデータと用語を定義し，\ref{sec_related}章で関連研究を整理する．
\ref{sec_label}章でラベルの設計の経緯を，\ref{sec_validation}章で検証データとその再現性を述べる．
\ref{sec_compliance}章で制約の順守の判定を，\ref{sec_text}章で VLM の教師と文の埋め込みの分類器による判定を述べる．
\ref{sec_vlm}章で小さな VLM の二段階の学習を，\ref{sec_outside}章で画像外の知識の改善の試みを述べ，\ref{sec_speed}章で推論の高速化を示す．
\ref{sec_discussion}章で得られた知見と限界を考察し，\ref{sec_conclusion}章でまとめる．
